\section{预训练优化}
\label{sec:nlp-llm-optimization}

训练分布、网络结构和学习目标确定以后，还要决定怎样在有限预算内让目标值下降。规模扩大以后，小模型上靠反复试错解决的选择都变得昂贵：一次完整训练可能占用大量加速器数周到数月，配置错误或中途发散都要付出同等代价。本节回答五个问题：

\begin{enumerate}
  \item 计算预算怎样在模型参数量与训练词元数之间分配；
  \item 在小模型上调好的学习率与初始化，能否按规则迁移到大模型；
  \item 用什么优化器把梯度变成参数更新；
  \item 批次规模与学习率怎样随训练进程安排；
  \item 低精度运算下怎样防止数值溢出、logit漂移和损失尖峰。
\end{enumerate}

对应的方法沿这五个问题展开。缩放律用小规模实验拟合损失与规模的经验关系，从Kaplan的幂律到Chinchilla的计算最优分配，再到计入推断成本后的超额训练；最大更新参数化让最优超参数不随宽度漂移，使窄模型上的调参结果可以直接放大；AdamW是默认优化器，Muon对矩阵更新做正交化，并通过尺度匹配沿用AdamW的超参数；预热加余弦衰减是经典日程，WSD日程去掉对总步数的依赖，临界批次规模给出增大批次的收益上限；稳定性措施分别作用于数值表示、梯度、输出logit和注意力logit。

本节用到四项一般结论：Adam用梯度的一阶矩与二阶矩逐坐标缩放更新；余弦调度在预定步数内把学习率平滑降到下限；全局范数裁剪在梯度过大时按比例缩小它而不改方向；Newton--Schulz矩阵迭代只用矩阵乘法就能把矩阵的非零奇异值推向1。它们的推导可参考\BookChapterRef{chap:neural-network-training}{卷五“神经网络模型”的“神经网络的训练”章}。这里只写大语言模型特有的配置、推导和比较；设备间通信与并行切分不在本节范围内。

\subsection{规模选择与缩放律}
\label{sec:nlp-llm-compute-allocation}

在固定计算预算下，扩大模型意味着能处理的训练词元变少；扩大数据又受模型容量限制。参数多而曝光不足，模型尚未学够；模型太小，再多词元也受表达能力约束。要在训练之前作出选择，需要一个预测“给定参数量与训练词元数时损失有多大”的经验关系。下面先估计计算量，再比较两组缩放律的形式、推导和分歧来源，最后说明计入推断成本时为何常把模型训练到计算最优点之后。

\subsubsection{训练计算量的估计}
\label{sec:nlp-llm-compute-estimate}

第\ref{sec:nlp-llm-decoder-cost}小节已经由矩阵乘法逐项数出：每个词元前向约$2N$次浮点运算，反向要同时求激活梯度与权重梯度，约$4N$次，训练$D$个词元的计算量约为式\eqref{eq:nlp-llm-train-flops}的$C\approx6ND$。下文的缩放律都以这一估计作为预算约束，并沿用其两条适用条件：注意力中不含参数的两步乘法随上下文长度增长，在$T=8192$的LLaMA-3-8B上约为参数项的14\%，在13.1万词元时反超参数项，长上下文阶段必须另计；稀疏专家模型中$N$取激活参数量而非总参数量。

按式\eqref{eq:nlp-llm-train-flops}，在约$1.5\times10^{13}$个词元上训练8B模型，计算量约为$6\times8\times10^{9}\times1.5\times10^{13}\approx7.2\times10^{23}$次浮点运算。这个数在下文比较不同规模分配时反复使用。

\subsubsection{Kaplan缩放律与Chinchilla缩放律}
\label{sec:nlp-llm-scaling-laws}

\term{缩放定律}（scaling law，下文简称缩放律）指损失随参数量、数据量或计算量按幂函数下降的经验规律。Kaplan等人在一族仅解码器Transformer上发现，只受参数量限制时损失近似为$L(N)=(N_c/N)^{\alpha_N}$，只受数据量限制时为$L(D)=(D_c/D)^{\alpha_D}$，拟合指数约为$\alpha_N\approx0.076$、$\alpha_D\approx0.095$；这里的$N$不含嵌入参数。两者合并为
\begin{equation}
 L(N,D)=\left[\left(\frac{N_c}{N}\right)^{\alpha_N/\alpha_D}+\frac{D_c}{D}\right]^{\alpha_D},
 \label{eq:nlp-llm-kaplan}
\end{equation}
$N_c$、$D_c$为拟合常数\scite{kaplan2020scaling}。他们由此得到的计算高效前沿是$N_{\mathrm{opt}}\propto C^{0.73}$、$D_{\mathrm{opt}}\propto C^{0.27}$：预算增加10倍时，模型扩大约5.4倍，词元只增加约1.9倍。这条建议推动了一批参数量很大、训练词元相对较少的模型。

Chinchilla工作重新做了这一问题，用三种方法估计计算最优分配，其中参数化方法拟合
\begin{equation}
 L(N,D)=E+\frac{A}{N^{\alpha}}+\frac{B}{D^{\beta}},
 \qquad A,B,\alpha,\beta>0.
 \label{eq:nlp-llm-scaling-loss}
\end{equation}
$E$是数据本身不可约的损失，第二项是有限模型带来的额外损失，第三项是有限训练词元带来的额外损失\scite{hoffmann2022chinchilla}。\term{计算最优}（compute-optimal）分配指在$C=6ND$约束下使$L$最小的$(N,D)$。代入$D=C/(6N)$，并注意$\partial D/\partial N=-D/N$，有
\begin{equation*}
 \frac{\mathrm dL}{\mathrm dN}
 =-\frac{\alpha A}{N^{\alpha+1}}+\frac{\beta B}{D^{\beta+1}}\cdot\frac{D}{N}
 =\frac1N\left(-\alpha AN^{-\alpha}+\beta BD^{-\beta}\right).
\end{equation*}
令其为零，再把$D=C/(6N)$代回，解出
\begin{equation}
 \alpha AN^{-\alpha}=\beta BD^{-\beta},\qquad
 N_{\mathrm{opt}}=G\left(\frac C6\right)^{\frac{\beta}{\alpha+\beta}},\quad
 D_{\mathrm{opt}}=G^{-1}\left(\frac C6\right)^{\frac{\alpha}{\alpha+\beta}},
 \label{eq:nlp-llm-scaling-allocation}
\end{equation}
其中$G=(\alpha A/(\beta B))^{1/(\alpha+\beta)}$。驻点就是最小点：令$u=\ln N$，目标写成$E+Ae^{-\alpha u}+B(C/6)^{-\beta}e^{\beta u}$，两个指数项都是$u$的严格凸函数，和也严格凸。第一式有直接的经济含义：$\alpha AN^{-\alpha}=-N\,\partial L/\partial N$是参数量增加1\%带来的损失下降，$\beta BD^{-\beta}$是词元数增加1\%带来的损失下降；两者各消耗1\%的预算，最优点要求两种投入的边际收益相等。

Chinchilla的三种方法给出的两个指数都接近0.5，即预算增加时$N$与$D$应大致同比增长，比值约为每参数20个词元。按此设计的Chinchilla（70B参数、1.4T词元）与Gopher（280B参数、300B词元）计算量相近，前者在多数下游任务上更好。原文参数化拟合报告的四舍五入系数彼此不完全自洽，Besiroglu等人用原始数据重新拟合，得到$E\approx1.82$、$A\approx482$、$B\approx2085$、$\alpha\approx0.348$、$\beta\approx0.366$\scite{besiroglu2024chinchilla}。用其未取整的系数代入式\eqref{eq:nlp-llm-scaling-loss}，Chinchilla与Gopher的拟合损失分别约为1.974和2.000；在$C=5.76\times10^{23}$时，式\eqref{eq:nlp-llm-scaling-allocation}给出约72B参数、1.3T词元，每参数约18个词元，与原文结论一致。

“每参数约20个词元”成立的条件需要一并记住：训练材料基本不重复，数据来自MassiveText及其词元化，模型是稠密仅解码器结构，学习率按余弦日程且周期与训练长度一致，拟合范围限于原文实验的规模。换数据、换词元化或换结构后，系数都要重新估计。

两组结论为何分歧？已有分析指向两个主要原因。其一是参数计数：Kaplan使用不含嵌入的参数量，而他们的模型偏小，嵌入在总参数中占比不低；Pearce与Song按总参数重新计数，并把拟合限制在小规模区间，就能从Chinchilla式的规律复现出接近0.73的指数\scite{pearce2024reconciling}。其二是学习率日程：Kaplan对不同训练长度使用同一个余弦周期，提前停止的运行在学习率尚未衰减时被评估，它们的损失被高估，数据的价值被低估；Chinchilla指出余弦周期须与训练长度匹配。Porian等人进一步校正了输出层计算量、预热长度与小规模下的优化器超参数，校正后Kaplan的实验设置也得到与Chinchilla一致的分配\scite{porian2024resolving}。

\subsubsection{超出计算最优点的训练}
\label{sec:nlp-llm-beyond-compute-optimal}

计算最优只考虑训练成本。模型部署后每生成或读取一个词元约需$2N$次运算，若模型一生要处理$D_{\mathrm{inf}}$个推断词元，总成本为$6ND+2ND_{\mathrm{inf}}$。现在固定目标损失$L^\star$，求总成本最小的$(N,D)$。对拉格朗日函数$N(6D+2D_{\mathrm{inf}})+\lambda(E+AN^{-\alpha}+BD^{-\beta}-L^\star)$分别对$N$和$D$求导并令其为零：
\begin{equation*}
 6D+2D_{\mathrm{inf}}=\lambda\alpha AN^{-\alpha-1},\qquad
 6N=\lambda\beta BD^{-\beta-1}.
\end{equation*}
两式相除，整理得
\begin{equation}
 \frac{\alpha AN^{-\alpha}}{\beta BD^{-\beta}}=1+\frac{D_{\mathrm{inf}}}{3D}.
 \label{eq:nlp-llm-inference-aware}
\end{equation}
$D_{\mathrm{inf}}=0$时回到式\eqref{eq:nlp-llm-scaling-allocation}的平衡条件。推断量越大，右端越大，最优点要求增加参数的边际收益高于增加词元的边际收益，即停在参数较少、词元较多的位置：用更多一次性的训练计算换取此后每次调用都更便宜\scite{sardana2024beyond}。

LLaMA-3-8B约用15T词元训练，每参数近1900个词元，是Chinchilla比例的九十多倍\scite{nlp-llm-llama3-2024}。用上面的重拟合系数作一次外推（仅为教学计算，拟合本身未覆盖这一区间）：同样$7.2\times10^{23}$的预算，计算最优点约为81B参数、1.5T词元，拟合损失约1.968；8B、15T的拟合损失约2.022，略高，但每词元推断运算只有前者的十分之一。

超额训练需要足够的不重复材料。Muennighoff等人在数据受限条件下发现，重复约4轮以内的收益接近新数据，继续重复的收益很快衰减到接近零\scite{muennighoff2023scaling}。因此高质量材料的总量也会限制向右移动多远。稀疏专家模型的损失同时依赖总参数与激活参数，式\eqref{eq:nlp-llm-scaling-loss}的系数不能直接套用，需要重新拟合。

\subsection{参数化与超参数迁移}
\label{sec:nlp-llm-parameter-scale}

在目标规模上逐一搜索学习率和初始化，代价接近训练多个大模型。若能在窄模型上调参，再按规则迁移到宽模型，就需要让激活和每步更新的尺度都不随宽度变化。下面先说明这是两项不同的要求，再写满足它们的参数化规则及迁移流程。

\subsubsection{初始化与更新尺度}
\label{sec:nlp-llm-init-update-scale}

先看前向尺度。设输入$\bm h$为$n$维，一行权重$\bm w$的元素零均值、互相独立、方差为$\sigma^2$，并独立于输入。给定$\bm h$时，线性输出满足$\operatorname{Var}_{\bm w}(\bm w^{\top}\bm h\mid\bm h)=\sigma^2\lVert\bm h\rVert_2^2$；若各输入分量的二阶矩均为$\nu$，再对输入平均，输出方差为$n\sigma^2\nu$。取$\sigma^2\propto1/n$，即方差与扇入成反比，输出尺度就不随宽度变化。这是初始化时刻的计算，训练后权重与输入不再独立。

再看更新尺度。一步更新$\Delta\bm W$使输出变化$\Delta\bm y=\Delta\bm W\bm h$。与初始权重不同，$\Delta\bm W$由梯度$\bm\delta\bm h^{\top}$构成，与输入$\bm h$相关。对Adam一类更新，每个元素的幅度约为学习率$\eta$，符号与$\delta_i h_j$一致，于是
\begin{equation*}
 \Delta y_i=\sum_{j=1}^{n}\Delta W_{i,j}h_j\approx\eta\operatorname{sign}(\delta_i)\sum_{j=1}^{n}\lvert h_j\rvert,
\end{equation*}
$n$个同号项相加，结果随$n$线性增长。初始化时的随机求和只按$\sqrt n$增长（中心极限定理的尺度），更新的相关求和却按$n$增长（大数定律的尺度）。两种尺度不同，所以初始化合适不能推出更新合适。以教学示例说明：宽度从256增至4096，即16倍，同一学习率使隐藏特征的单步变化约增大16倍；要保持变化同阶，隐藏层学习率须按$1/n$缩小。

通常的初始化方差取$1/\mathrm{fan\_in}$、各层共用一个学习率，这种做法称为\term{标准参数化}（standard parametrization, SP）。在标准参数化下，最优学习率随宽度向小值漂移，小模型上调好的学习率放到大模型上偏大，甚至发散。

\subsubsection{\texorpdfstring{$\mu$P与$\mu$Transfer}{μP与μTransfer}}
\label{sec:nlp-llm-mup}

\term{最大更新参数化}（maximal update parametrization, $\mu$P）按参数所在位置分别规定初始化方差、学习率和前向乘子，使宽度增大时每层特征的更新都保持同阶，既不随宽度爆炸，也不消失\scite{nlp-llm-mutransfer2022}。表~\ref{tab:nlp-llm-mup}列出Adam优化器下的规则，宽度记为$n$，各量只写随$n$变化的比例。

\begin{table}[htbp]
 \centering
 \small
 \begin{tabular}{@{}llll@{}}
  \toprule
  参数位置 & 初始化方差 & Adam学习率 & 前向乘子\\
  \midrule
  输入层（嵌入）与偏置 & $\propto1$ & $\propto1$ & 1\\
  隐藏层矩阵 & $\propto1/n$ & $\propto1/n$ & 1\\
  输出层矩阵 & $\propto1/n$（可取零） & $\propto1$ & $\propto1/n$\\
  注意力打分 & — & — & $1/d_h$代替$1/\sqrt{d_h}$\\
  \midrule
  标准参数化（对照） & 均为$1/\mathrm{fan\_in}$ & 全局同一值 & 1，打分用$1/\sqrt{d_h}$\\
  \bottomrule
 \end{tabular}
 \caption{$\mu$P在Adam下的宽度缩放规则。比例相对于基准宽度，如宽度加倍时隐藏层学习率减半。输出层乘子与隐藏层学习率出于同一理由：二者处理的都是$n$个相关量的求和。}
 \label{tab:nlp-llm-mup}
\end{table}

表中各行都由第\ref{sec:nlp-llm-init-update-scale}小节的两种尺度得出。嵌入层的扇入是词表而不是宽度，按宽度不需缩放。隐藏层的更新是$n$项相关求和，学习率乘$1/n$。输出层把$n$维特征汇总成logit，训练后特征与输出权重相关，求和按$n$增长，因此在前向乘以$1/n$；有了这个乘子，输出层的学习率就不必再缩放。注意力中训练后的$\bm q$与$\bm k$也会相关，点积随$d_h$线性增长，$\mu$P用$1/d_h$代替$1/\sqrt{d_h}$。

$\mu$P的直接用途是$\mu$Transfer：最优超参数在$\mu$P下近似不随宽度变化，于是可以在窄代理模型上调参，再把结果原样放到宽模型上，流程见算法~\ref{alg:nlp-llm-mutransfer}。

\begin{algorithm}[$\mu$Transfer超参数迁移]
  \label{alg:nlp-llm-mutransfer}
  \begin{algorithmic}[1]
    \Require 目标宽度$n$的模型结构，基准宽度$n_0\ll n$，待调超参数集合$\mathcal H$（学习率、初始化尺度、输出乘子等）
    \Ensure 目标模型的超参数
    \State 按表~\ref{tab:nlp-llm-mup}以$n_0$为基准写出$\mu$P参数化，使宽度为$n_0$时与标准参数化一致
    \State 构造宽度为$n_0$、其余结构与目标相同的代理模型
    \For{$\mathcal H$中的每组候选}
      \State 用较少词元训练代理模型，记录留出损失
    \EndFor
    \State 取留出损失最低的一组$h^\star$
    \State 在宽度$n$上按$\mu$P规则换算各参数组的初始化、学习率与乘子，其余直接使用$h^\star$
    \State 训练目标模型，并在早期检查激活与更新尺度是否与代理模型同阶
  \end{algorithmic}
\end{algorithm}

Yang等人用约4000万参数的代理调参，迁移到67亿参数的GPT-3结构，结果优于原文公布的同规模模型，调参成本约为一次预训练的7\%\scite{nlp-llm-mutransfer2022}。$\mu$Transfer的理论保证只针对宽度；深度、批次、训练长度的迁移来自经验观察，正则化强度等依赖数据的超参数也不一定可迁移。它决定的是给定规模下怎样设置超参数，与式\eqref{eq:nlp-llm-scaling-allocation}决定选多大规模是两个不同问题，二者可以配合：先在$\mu$P下做小规模缩放律实验，再外推到目标规模。

\subsection{优化器}
\label{sec:nlp-llm-optimizer}

大语言模型训练的默认优化器是AdamW，它对每个坐标单独估计梯度尺度。Muon把二维权重矩阵当作整体，用正交化后的动量作为更新方向。下面比较这两种思路，并说明在规模化训练中怎样让两者的更新尺度一致。

\subsubsection{AdamW}
\label{sec:nlp-llm-adamw-update}

采用\term{解耦权重衰减}（decoupled weight decay）的Adam（Adam with decoupled weight decay, AdamW）维护一阶矩$\bm m$和二阶矩$\bm v$。记第$t$步梯度为$\bm g^{(t)}$，更新为
\begin{align}
 \bm m^{(t)}&=\beta_1\bm m^{(t-1)}+(1-\beta_1)\bm g^{(t)},\qquad
 \bm v^{(t)}=\beta_2\bm v^{(t-1)}+(1-\beta_2)(\bm g^{(t)})^{\odot2},\notag\\
 \bm\theta^{(t+1)}&=\bm\theta^{(t)}
  -\eta_t\left(\frac{\widehat{\bm m}^{(t)}}{\sqrt{\widehat{\bm v}^{(t)}}+\varepsilon}+\lambda\bm\theta^{(t)}\right).
 \label{eq:nlp-llm-adamw}
\end{align}
其中$\beta_1,\beta_2\in[0,1)$为平滑系数；$\widehat{\bm m}^{(t)}=\bm m^{(t)}/(1-\beta_1^t)$、$\widehat{\bm v}^{(t)}=\bm v^{(t)}/(1-\beta_2^t)$校正零初始化带来的偏差；除法、平方根和$(\cdot)^{\odot2}$逐坐标进行；$\eta_t$为学习率，$\lambda\geq0$为权重衰减系数，$\varepsilon>0$防止分母为零\scite{loshchilov2019adamw}。

“解耦”是相对L2正则而言的。若把$\tfrac\lambda2\lVert\bm\theta\rVert_2^2$加进损失，梯度变为$\bm g+\lambda\bm\theta$，衰减项也要除以$\sqrt{\widehat v_i}+\varepsilon$；于是历史梯度大的坐标衰减弱，历史梯度小的坐标衰减强，衰减强度被自适应缩放扭曲。式\eqref{eq:nlp-llm-adamw}把$\lambda\bm\theta$放在预处理之外，每个坐标按同一比例$1-\eta_t\lambda$收缩。对SGD两者等价，对Adam不等价。

大语言模型训练常用$\beta_1=0.9$、$\beta_2=0.95$、$\lambda=0.1$，归一化层的增益和偏置不加衰减。$\beta_2$比通用默认值0.999小，二阶矩能更快跟上梯度尺度的变化；这是工程惯例。AdamW为每个参数保存$m_i$和$v_i$两个状态，通常用fp32存储，优化器状态占参数量两倍的存储；第\ref{chap:nlp-supervised-finetuning}章估算微调显存时会用到这一点。

\subsubsection{Muon与矩阵正交化更新}
\label{sec:nlp-llm-muon}

AdamW逐坐标缩放，看不到矩阵的行列结构。对二维权重，动量矩阵的奇异值往往很不均衡：少数方向的奇异值很大，更新几乎只沿这几个方向进行。Muon的思路是保留奇异向量、把奇异值统一为1，使各方向以相同幅度更新\scite{jordan2024muon}。对隐藏层矩阵$\bm W\in\mathbb R^{m\times n}$及其梯度$\bm G^{(t)}$，Muon计算
\begin{equation}
 \begin{aligned}
 \bm M^{(t)}&=\beta_1\bm M^{(t-1)}+\bm G^{(t)},\qquad
 \bm O^{(t)}=\operatorname{NS}_5\bigl(\bm M^{(t)}\bigr)\approx\bm U\bm V^{\top},\\
 \bm W^{(t+1)}&=\bm W^{(t)}-\eta_t\bm O^{(t)},
 \end{aligned}
 \label{eq:nlp-llm-muon}
\end{equation}
其中$\bm M=\bm U\bm\Sigma\bm V^{\top}$为紧奇异值分解，$\beta_1$为动量系数。$\operatorname{NS}_5$先令$\bm X_0=\bm M/\lVert\bm M\rVert_{\mathrm F}$，再执行5次Newton--Schulz迭代$\bm X\gets a\bm X+b(\bm X\bm X^{\top})\bm X+c(\bm X\bm X^{\top})^2\bm X$。每次迭代对每个奇异值作用同一个五次多项式$\sigma\mapsto a\sigma+b\sigma^3+c\sigma^5$，原实现取$(a,b,c)=(3.4445,-4.7750,2.0315)$，使$(0,1]$内的奇异值在5步内都被推到1附近。这组系数追求收敛快，结果并不精确等于1，原作者报告这种近似不影响训练效果；整个过程只需矩阵乘法，不做奇异值分解。嵌入层、输出层和增益、偏置等非矩阵参数仍用AdamW。

Muon用于大语言模型时需要处理更新尺度。半正交矩阵$\bm O$的Frobenius范数平方等于$\min(m,n)$，元素均方根为
\begin{equation*}
 \operatorname{rms}(\bm O)=\sqrt{\frac{\min(m,n)}{mn}}=\frac{1}{\sqrt{\max(m,n)}}.
\end{equation*}
形状不同的矩阵，更新的元素尺度不同；它与AdamW的更新尺度也不同，后者的元素均方根经验上约为0.2到0.4。Moonlight的做法是把更新改为$0.2\sqrt{\max(m,n)}\,\bm O^{(t)}$，使各矩阵的更新均方根都为0.2，与AdamW匹配，从而可以沿用AdamW的学习率与权重衰减；同时加入解耦权重衰减，防止长时间训练中权重和输出幅度持续增长\scite{nlp-llm-muonscale2025}。在他们的缩放律实验中，达到同一损失所需的计算约为AdamW的一半。

两种优化器的差别落在三处。更新方向上，AdamW按各坐标的历史平方梯度重标度，Muon按矩阵的奇异方向重标度；被正交化的是更新，权重本身并不被约束为正交矩阵。存储上，Muon每个矩阵只存一份动量，比AdamW少一份状态。计算上，每步多出5次矩阵迭代，$m\times n$矩阵每次约$O(mn\min(m,n))$次运算，在大规模分布式训练中还要先汇集完整矩阵。Muon的收益报告来自特定结构与数据，替换优化器时仍需在相同预算下重新比较。

\subsection{批次与学习率日程}
\label{sec:nlp-llm-training-schedule}

学习率决定每步走多远，批次决定梯度估计的噪声，二者要一起安排。先确定批次的计量口径：若每个目标词元等权，一批含$B$个计分词元时，梯度是这些词元负对数似然梯度的平均。在样本独立、方差有限的理想情况下，平均梯度的方差按$1/B$下降；但长文档中相邻位置并不独立，重复样本也不提供新信息，计分词元数不能直接当作同样数量的独立观测。

日程的横轴也要说清。设两个配置都训练一百万步，第二个每步的计分词元数是第一个的4倍，它们的数据曝光和计算量都相差4倍。若批次在训练中逐步增大，而学习率仍按步数衰减，每个学习率阶段处理的词元数就随之改变。比较时应同时记录步数、累计词元数、学习率与留出损失。

\subsubsection{预热与余弦衰减}
\label{sec:nlp-llm-cosine}

最常见的日程是先线性预热，再余弦衰减。设总步数为$S$，预热步数为$S_{\mathrm w}$，峰值与下限学习率为$\eta_{\max}$、$\eta_{\min}$：
\begin{equation}
 \eta_t=
 \begin{cases}
  \eta_{\max}\,t/S_{\mathrm w}, & t\leq S_{\mathrm w},\\[2pt]
  \eta_{\min}+\tfrac12(\eta_{\max}-\eta_{\min})\left(1+\cos\dfrac{\pi(t-S_{\mathrm w})}{S-S_{\mathrm w}}\right), & S_{\mathrm w}<t\leq S.
 \end{cases}
 \label{eq:nlp-llm-cosine}
\end{equation}
预热的理由有两条：训练初期Adam的二阶矩只由少数几步梯度估计，方差大，用它缩放的步长不可靠\scite{liu2020radam}；后置归一化等结构在初期梯度很大，需要小步长起步\scite{xiong2020layernorm}。LLaMA取2000步预热，余弦衰减到峰值的10\%\scite{touvron2023llama}。

余弦日程的缺点是总步数$S$必须事先确定。训练到一半若想延长，学习率已经下降，无法简单接着训练；而Chinchilla的分析表明，余弦周期与实际训练长度不匹配会使最终损失变差。这使得“在多个训练长度上各得一个模型”需要多次从头训练。

\subsubsection{WSD日程与退火}
\label{sec:nlp-llm-wsd}

\term{预热—稳定—衰减日程}（warmup-stable-decay schedule, WSD）把学习率分为三段：预热到$\eta_{\max}$，保持恒定，最后在较短的衰减段降到接近零\scite{hu2024minicpm}。设恒定段结束于$S_{\mathrm s}$、衰减段长$S_{\mathrm d}$，衰减段一种常用形式为
\begin{equation}
 \eta_t=\eta_{\max}\,f\!\left(\frac{t-S_{\mathrm s}}{S_{\mathrm d}}\right),\qquad
 S_{\mathrm s}<t\leq S_{\mathrm s}+S_{\mathrm d},
 \label{eq:nlp-llm-wsd}
\end{equation}
其中$f(0)=1$、$f(1)\approx0$，可取线性$f(u)=1-u$或$f(u)=1-\sqrt u$；衰减段通常占总步数的10\%到20\%。

图~\ref{fig:nlp-llm-lr-schedules}对比两种日程。WSD的关键性质是恒定段与总长度无关：恒定段中的任何检查点都可以分出一个短衰减段，得到“训练到此为止”的完整模型，主干则继续以$\eta_{\max}$训练。一次主干运行因此能得到多个训练长度的终态模型，缩放律实验和后续继续训练都不必从头开始。Hägele等人比较发现，衰减段占比约20\%时，WSD的终态损失与按长度调好的余弦日程相当\scite{hagele2024scaling}。

\begin{figure}[htbp]
  \centering
  % 预热加余弦衰减与WSD日程的学习率曲线（解析式绘制，参数为教学示例）。
% 余弦：前5%线性预热，之后按余弦降到0.1；WSD：前5%预热，恒定到80%，最后20%线性降到0；
% 虚线为从40%与60%处恒定段检查点分出的短衰减段，长度取各自已训练步数的20%。
\begin{tikzpicture}[figure picture]
  \begin{axis}[
    width=\linewidth,
    height=54mm,
    xmin=0,xmax=100,
    ymin=0,ymax=1.2,
    xlabel={训练进度（占计划总步数的百分比）},
    ylabel={$\eta_t/\eta_{\max}$},
    xtick={0,20,40,60,80,100},
    ytick={0,0.5,1},
    grid=major,
    legend style={at={(0.5,1.03)},anchor=south,legend columns=3,
      /tikz/every even column/.append style={column sep=4mm}},
  ]
    \addplot[FigureBlue,line width=1.1pt,domain=0:5,samples=2,forget plot] {x/5};
    \addplot[FigureBlue,line width=1.1pt,domain=5:100,samples=120]
      {0.1+0.45*(1+cos(180*(x-5)/95))};
    \addlegendentry{预热加余弦}
    \addplot[FigureRed,line width=1.1pt,domain=0:5,samples=2,forget plot] {x/5};
    \addplot[FigureRed,line width=1.1pt,domain=5:80,samples=2,forget plot] {1};
    \addplot[FigureRed,line width=1.1pt,domain=80:100,samples=2] {(100-x)/20};
    \addlegendentry{WSD}
    \addplot[FigureYellow!80!black,line width=1pt,dashed,domain=40:48,samples=2] {(48-x)/8};
    \addlegendentry{分支衰减}
    \addplot[FigureYellow!80!black,line width=1pt,dashed,domain=60:72,samples=2,forget plot] {(72-x)/12};
    \addplot[only marks,mark=*,mark size=1.4pt,FigureStroke,forget plot] coordinates {(40,1) (60,1)};
    \node[anchor=south,font=\figurefont\footnotesize,text=FigureStroke] at (axis cs:50,1.01) {分支检查点};
  \end{axis}
\end{tikzpicture}

  \caption{两种学习率日程的形状。横轴为占计划总步数的百分比，纵轴为相对峰值的学习率；参数为教学示例：预热5\%，余弦降到峰值的10\%，WSD在最后20\%线性降到零。虚线表示从恒定段40\%与60\%处的检查点分出的短衰减段，长度为已训练步数的20\%；余弦日程没有可以这样分支的恒定段。}
  \label{fig:nlp-llm-lr-schedules}
\end{figure}

衰减段中训练损失快速下降。一种直观解释是：大学习率下，参数在损失曲面陡峭方向上被梯度噪声持续推动，停在离谷底一定距离处振荡，这段距离随学习率增大；学习率降低后振荡减小，损失落向谷底。衰减段称为\term{学习率退火}（learning rate annealing）。由于这一段对最终模型影响大，实践中常在退火时提高高质量数据（如数学、代码和精选文本）的比例。LLaMA 3在预训练末尾以高质量数据退火，并对退火阶段的检查点取平均\scite{nlp-llm-llama3-2024}。这种安排与第\ref{sec:nlp-llm-continuation}节的继续训练、第\ref{chap:nlp-reasoning}章推理能力的预训练基础相衔接。

\subsubsection{临界批次规模}
\label{sec:nlp-llm-critical-batch}

增大批次可以减少步数，但收益有上限。McCandlish等人用二次模型给出了这个上限\scite{mccandlish2018empirical}。设真实梯度为$\bm G$，单样本梯度的协方差为$\bm\Sigma$，批次为$B$时批次梯度$\bm g$满足$\mathbb E\bm g=\bm G$、$\operatorname{Cov}\bm g=\bm\Sigma/B$。损失在当前点的局部曲率为$\bm H$，沿$-\epsilon\bm g$走一步，损失的期望变化为
\begin{equation*}
 \mathbb E[\Delta L]\approx-\epsilon\lVert\bm G\rVert_2^2+\tfrac12\epsilon^2\left(\bm G^{\top}\bm H\bm G+\frac{\operatorname{tr}(\bm H\bm\Sigma)}{B}\right).
\end{equation*}
对$\epsilon$求最优，得到
\begin{equation}
 \epsilon_{\mathrm{opt}}(B)=\frac{\epsilon_{\max}}{1+B_{\mathrm{noise}}/B},\qquad
 \Delta L_{\mathrm{opt}}(B)=\frac{\Delta L_{\max}}{1+B_{\mathrm{noise}}/B},\qquad
 B_{\mathrm{noise}}=\frac{\operatorname{tr}(\bm H\bm\Sigma)}{\bm G^{\top}\bm H\bm G},
 \label{eq:nlp-llm-noise-scale}
\end{equation}
其中$\epsilon_{\max}=\lVert\bm G\rVert_2^2/(\bm G^{\top}\bm H\bm G)$和$\Delta L_{\max}$是用无噪声梯度时的最优步长与最大下降。若$\bm H$近似各向同性，$B_{\mathrm{noise}}$化为$B_{\mathrm{simple}}=\operatorname{tr}(\bm\Sigma)/\lVert\bm G\rVert_2^2$，即梯度噪声的总方差与梯度信号强度之比，称为\term{梯度噪声尺度}（gradient noise scale）。它可以在训练中用不同批次的梯度范数估计。

式\eqref{eq:nlp-llm-noise-scale}说明，$B\ll B_{\mathrm{noise}}$时每步下降近似与$B$成正比，批次加倍几乎让所需步数减半，处理的词元总数不变；$B\gg B_{\mathrm{noise}}$时每步下降接近上限，再增大批次只是多消耗数据。两区之间的过渡点称为\term{临界批次规模}（critical batch size），与$B_{\mathrm{noise}}$同阶。以教学示例说明：设$B_{\mathrm{noise}}$为100万词元，批次为25万、100万、400万词元时，单步下降分别为$\Delta L_{\max}$的0.2、0.5、0.8倍。达到同一进展所需步数之比为$5:2:1.25$，所需词元数之比为$1.25:2:5$。从25万增至100万，步数减少到原来的40\%，词元只多用60\%；从100万增至400万，步数只减少到62.5\%，词元却多用1.5倍。

同一公式还给出两条实践含义。第一，$\epsilon_{\mathrm{opt}}$随$B$增大，批次与学习率需要联动调整，不能只改一个。第二，训练进行中$\lVert\bm G\rVert_2$变小，$B_{\mathrm{noise}}$增大；Kaplan等人观察到临界批次规模随损失下降而增大\scite{kaplan2020scaling}。这解释了批次递增日程：GPT-3在前40亿到120亿个词元内把批次从3.2万词元线性增大到最终值\scite{brown2020}。式\eqref{eq:nlp-llm-noise-scale}基于局部二次近似和最优步长的假设，给出的是量级判断。

\subsection{训练稳定性}
\label{sec:nlp-llm-stability}

大规模训练常出现损失突然升高、梯度范数暴涨甚至发散。来源大致有三类：低精度表示的溢出与舍入，输出logit和注意力logit的持续增长，以及个别数据批次与当前参数状态的异常组合。下面按“数值表示—梯度与输出logit—注意力logit与事后处理”的顺序给出措施及其作用对象。DeepSeek-V3等技术报告同时采用多项措施\scite{deepseek2024v3}，单份报告不能分离每项措施的独立贡献。

\subsubsection{混合精度训练}
\label{sec:nlp-llm-mixed-precision}

\term{混合精度训练}（mixed-precision training）用低精度格式做前向与反向的矩阵乘法，用fp32保存需要精细累积的量\scite{micikevicius2018mixed}。fp16有5位指数、10位尾数，最大值约65504，最小正规数约$6\times10^{-5}$，梯度容易下溢，需要先把损失乘一个大系数再反算（损失缩放）。bf16有8位指数、7位尾数，指数位与fp32相同，动态范围一致，不需要损失缩放；代价是精度低，1附近相邻两个可表示数相差$2^{-7}\approx0.0078$。

精度不足会吞掉小更新。设某个权重为1.0，一步更新为$10^{-4}$，在bf16中$1.0+10^{-4}$舍入后仍为1.0，这步更新完全丢失；成千上万步的小更新都会这样丢失。因此参数要另存一份fp32的\term{主权重}（master weights）：前向与反向使用bf16副本，优化器在fp32主权重上累积更新，再转换出新的bf16副本。按每个参数计，bf16权重与梯度各2字节，fp32主权重4字节，Adam的两个矩各4字节，合计约16字节；8B模型仅这几项就需约128GB，还不含激活。

更进一步的FP8训练把矩阵乘法降到8位浮点。FP8动态范围很窄，需要按小块分别计算缩放系数，例如DeepSeek-V3对激活按$1\times128$、对权重按$128\times128$的块缩放，并在更高精度中累加\scite{deepseek2024v3}。

\subsubsection{梯度裁剪与z-loss}
\label{sec:nlp-llm-zloss}

全局范数裁剪把所有参数的梯度拼成一个向量$\bm g$，按$\bm g\gets\bm g\cdot\min(1,c/\lVert\bm g\rVert_2)$处理，阈值$c$常取1.0。它保留方向，只限制单个异常批次造成的更新幅度。经过Adam的逐坐标缩放后，梯度范数上限并不等于参数步长上限；若几乎每步都被裁剪，说明梯度尺度持续偏大，应检查学习率或初始化。

输出logit还有一个裁剪管不到的问题。softmax对所有logit加同一常数不变，因此交叉熵约束不了logit的整体偏移，归一化项$\log Z=\log\sum_{v}\exp z_v$可以在训练中持续漂移。$\log Z$很大时，logit本身很大，bf16下的舍入误差随之放大，softmax可能饱和。z-loss在每个位置的损失中加入
\begin{equation}
 \ell_z=\lambda_z(\log Z)^2,\qquad
 \frac{\partial\ell_z}{\partial z_v}=2\lambda_z\log Z\cdot p_v,
 \label{eq:nlp-llm-zloss}
\end{equation}
$p_v$为softmax概率。梯度把所有logit朝使$\log Z$接近零的方向平移，而不改变它们之间的差，因此不影响预测分布本身。PaLM取$\lambda_z=10^{-4}$并报告它改善了训练稳定性\scite{chowdhery2023palm}。稀疏专家的路由器z-loss（式\eqref{eq:nlp-llm-router-z-loss}）形式相同，约束的是路由logit。

\subsubsection{QK归一化与损失尖峰处理}
\label{sec:nlp-llm-qknorm-spikes}

注意力打分$s=\bm q^{\top}\bm k/\sqrt{d_h}$没有上界。若训练中查询与键的范数持续增长，某些打分远大于其他打分，softmax退化为接近独热的分布，注意力熵快速下降，这种现象称为\term{注意力熵塌缩}（attention entropy collapse）\scite{zhai2023entropy}。Dehghani等人在220亿参数视觉Transformer中观察到注意力logit增长导致发散\scite{dehghani2023vit22b}；Wortsman等人发现，同类不稳定在小模型配合大学习率时就能复现，并可用小模型检验对策\scite{wortsman2024proxies}。

\term{查询--键归一化}（QKNorm）在打分前分别对每个头的$\bm q$和$\bm k$做RMSNorm或LayerNorm\scite{henry2020qknorm}。增益取1时，RMSNorm输出的范数恰为$\sqrt{d_h}$，由柯西--施瓦茨不等式，$\lvert\hat{\bm q}^{\top}\hat{\bm k}\rvert/\sqrt{d_h}\leq\lVert\hat{\bm q}\rVert_2\lVert\hat{\bm k}\rVert_2/\sqrt{d_h}=\sqrt{d_h}$。$d_h=128$时打分绝对值不超过约11.3，可学习增益再按需要放大，但增长受到直接约束。QK归一化要求每个头的键被显式算出；多头潜在注意力把键压缩在共享潜在向量中，不便直接插入归一化。Kimi K2改为在打分最大值超过阈值时缩小查询与键的投影权重\scite{kimiteam2025k2}。

即使采用以上措施，损失尖峰仍会出现。\term{损失尖峰}（loss spike）指损失在少数步内突然升高，可能自行回落，也可能导致发散。PaLM报告，从尖峰前约100步的检查点重启，并跳过约200到500个批次，尖峰就不再出现；用不同检查点重放同一批数据也未必复现尖峰，说明尖峰来自特定数据与特定参数状态的组合，而不是单纯的坏数据\scite{chowdhery2023palm}。常用的处理流程见算法~\ref{alg:nlp-llm-spike-recovery}。

\Needspace{16\baselineskip}
\begin{algorithm}[损失尖峰的监控与回滚]
  \label{alg:nlp-llm-spike-recovery}
  \begin{algorithmic}[1]
    \Require 定期保存的检查点，滑动窗口内的损失均值$\bar\ell$与标准差$s_\ell$，阈值$k$，回退步数$r$
    \Ensure 继续训练的参数与数据位置
    \For{每个训练步$t$}
      \State 记录损失$\ell_t$、梯度范数、各层更新与参数范数之比、最大注意力打分与$\log Z$
      \If{$\ell_t>\bar\ell+k\,s_\ell$且随后数步未回落}
        \State 载入第$t-r$步附近的检查点及其优化器状态
        \State 跳过第$t-r$步至第$t$步对应的数据批次
        \If{同一位置附近反复出现尖峰}
          \State 降低学习率，或检查上述监控量中最先异常的一项
        \EndIf
      \EndIf
    \EndFor
  \end{algorithmic}
\end{algorithm}

监控量帮助定位尖峰在哪一层、由哪类量最先失控，但单次尖峰通常无法确定唯一原因。回滚与跳批的代价是重算一段训练并丢弃少量数据，保存检查点的频率决定这段代价有多大。

表~\ref{tab:nlp-llm-optimization-methods}汇总本节方法。规模选择决定训练什么、训练多久；参数化与优化器决定每步怎样更新；日程决定更新幅度怎样随时间变化；稳定性措施则保证这一过程在有限精度下不中断。预算终点得到的参数是否适合交给下一阶段，要由第\ref{sec:nlp-llm-continuation}节的检查回答。

\begin{table*}[htbp]
 \small
 \begin{tabularx}{\linewidth}{@{}p{28mm}p{38mm}p{44mm}X@{}}
  \toprule
  方法 & 解决的问题 & 核心机制 & 新代价或适用边界\\
  \midrule
  Kaplan缩放律 & 预测损失随规模的变化 & 不含嵌入参数的幂律拟合 & 计数与日程设置使数据价值被低估\\
  Chinchilla缩放律 & 固定计算下分配$N$与$D$ & 拟合$E+A/N^\alpha+B/D^\beta$，边际收益相等 & 依赖数据、结构与日程，不计推断成本\\
  超额训练 & 降低部署后的推断成本 & 推断量大时偏向小模型、多词元 & 需要足够不重复的数据\\
  $\mu$P与$\mu$Transfer & 超参数随宽度漂移 & 按位置缩放初始化、学习率与乘子 & 理论只覆盖宽度；需改写参数化\\
  AdamW & 梯度尺度各坐标不同 & 逐坐标二阶矩缩放，衰减与预处理分离 & 两份优化器状态\\
  Muon & 矩阵更新集中在少数方向 & Newton--Schulz正交化动量，按形状匹配尺度 & 额外矩阵乘法；非矩阵参数另用AdamW\\
  预热加余弦 & 初期统计不准，后期需细调 & 线性升至峰值后余弦衰减 & 需预知总步数\\
  WSD日程 & 训练长度不确定 & 恒定段可随时分出短衰减段 & 衰减段长度与形状需另调\\
  临界批次规模 & 批次增大的收益上限 & 梯度噪声尺度$\operatorname{tr}\bm\Sigma/\lVert\bm G\rVert_2^2$ & 二次近似下的量级判断\\
  bf16与主权重 & 低精度溢出与小更新丢失 & bf16计算，fp32累积更新 & 每参数约16字节状态\\
  梯度裁剪与z-loss & 异常梯度与输出logit漂移 & 限制梯度范数；惩罚$(\log Z)^2$ & 裁剪不限制最终步长\\
  QK归一化 & 注意力logit增长与熵塌缩 & 打分前归一化查询与键 & 与潜在压缩的键不易兼容\\
  回滚与跳批 & 损失尖峰 & 回到早前检查点并跳过对应数据 & 重算开销，原因难以单独确定\\
  \bottomrule
 \end{tabularx}
 \caption{预训练优化方法对照。“新代价或适用边界”一列只列与本节讨论相关的主要项；各方法可以同时使用，表中顺序对应本节小节顺序。}
 \label{tab:nlp-llm-optimization-methods}
\end{table*}
